\subsubsection{循证阶梯与常见误区}

早泄治疗应遵循\textbf{按需、按型、按困扰程度}的原则，主流指南的阶梯大致是：\textbf{一线}为行为疗法（挤压法、停-动法、性感集中训练）与按需药物（达泊西汀、利多卡因/丙胺卡因局部麻醉剂）并行；\textbf{二线}为每日服用的 SSRI（多为超说明书）以及合并勃起功能障碍时联用 PDE5 抑制剂，行为与药物联合常优于单一手段；\textbf{三线}的\textbf{选择性阴茎背神经切断术已被主流指南明确不推荐}——缺乏可靠的长期获益证据，且可致永久性龟头感觉减退甚至勃起功能障碍，切勿轻信"手术断根"的宣传。

需要特别澄清的几类误区：

- \textbf{"时间越长越好"}：性满意度与时长并非线性关系，执着于"计时"反而加重表现焦虑与早泄。
- \textbf{自行购买"延时喷剂/湿巾"}：此类产品多含利多卡因、苯佐卡因等局麻成分，浓度剂量不明，过量可致龟头麻木、影响伴侣感觉，甚至不射精或勃起不坚。
- \textbf{"纯中药壮阳可根治"}：多数缺乏循证依据，长期滥用可能扰乱内分泌、损伤肝肾。
- \textbf{把"偶尔过快"当成病}：情境性的偶发射精过快属于正常波动（自然变异性早泄），不必用药；反之，把"勃起维持困难时的抢时间"误当早泄（早泄样射精功能障碍）也很常见，其处理重点是信心与认知，而非加药。

\subsubsection{伴侣参与与心理关系因素}

早泄几乎总在关系中出现、也在关系中缓解。伴侣的理解与配合能显著提高行为疗法的效果：把目标从"坚持多少分钟"改为"双方是否满意"；前戏充分、节奏协商、借助体位与手法加以补充；挤压法、停-动法最好由伴侣协助进行，并事先约定"暂停信号"。若存在表现焦虑、抑郁、关系积怨或性创伤，应寻求认知行为治疗或专业性与关系治疗；合并勃起功能障碍、慢性前列腺炎、甲状腺功能亢进或药物副作用者，须先治疗原发病。
